% =============================================================================
% CHAPTER 5 — RESEARCH METHODOLOGY
% Thesis: PATHCAST
% Covers: DSR, dataset, experimental design, metrics, baselines, threats
% Target: ~12-16 pages
% =============================================================================

\chapter{Research Methodology}
\label{ch:methodology}

This chapter describes the methodological framework adopted in this 
research. Section~\ref{sec:dsr} presents the Design Science Research 
paradigm and its instantiation. Section~\ref{sec:dataset} describes the 
dataset, including repository selection criteria and descriptive 
statistics. Section~\ref{sec:eval-design} formalizes the experimental 
design for the two evaluation stages. Section~\ref{sec:evaluation-metrics} 
defines the evaluation metrics for each research question. 
Section~\ref{sec:baselines} specifies the baseline methods against which 
the proposed approach is compared. Section~\ref{sec:statistical-analysis} 
describes the statistical analysis procedures. 
Section~\ref{sec:threats} discusses threats to validity and the 
mitigation strategies adopted.


% =============================================================================
% 5.1 DESIGN SCIENCE RESEARCH
% =============================================================================
\section{Design Science Research Paradigm}
\label{sec:dsr}

This research adopts the Design Science Research (DSR) paradigm, following 
the framework of Hevner et al.~\cite{hevner2004} and the Design 
Science Research Methodology (DSRM) process of Peffers et 
al.~\cite{peffers2007}. DSR is appropriate because the primary goal 
is to design, implement, and evaluate an artifact (the PATHCAST 
method and its supporting tool) that addresses a practical problem 
identified in the literature.


\subsection{DSR Artifact Characterization}
\label{sec:dsr-artifact}

According to the taxonomy of Hevner et al.~\cite{hevner2004}, DSR 
artifacts can be classified as constructs, models, methods, or 
instantiations. The artifact produced in this research is a \textbf{method}: 
PATHCAST is a formally specified pipeline that prescribes how to 
transform software repository data into probabilistic process forecasts 
through a defined sequence of operations.

The method artifact is accompanied by a supporting \textbf{instantiation}: 
a Python-based tool that implements the four pipeline stages and the ML 
component, serving as the vehicle for empirical evaluation.


\subsection{DSRM Process Instantiation}
\label{sec:dsrm-instantiation}

Table~\ref{tab:dsrm-mapping} maps the six DSRM activities to the 
corresponding thesis chapters and deliverables.

\begin{table}[htbp]
\centering
\caption{Mapping of DSRM activities to thesis components.}
\label{tab:dsrm-mapping}
\small
\begin{tabular}{clll}
\toprule
\textbf{Activity} & \textbf{DSRM Activity} & \textbf{Chapter} & \textbf{Deliverable} \\
\midrule
1 & Problem Identification & Ch.~\ref{ch:introduction} & 
    Three research problems \\
  & \& Motivation & Ch.~\ref{ch:theoretical-foundation} & 
    Gap analysis table \\
\midrule
2 & Objectives Definition & Ch.~\ref{ch:introduction} & 
    3 RQs, 4 specific objectives \\
  & & & 5 expected contributions \\
\midrule
3 & Design \& Development & Ch.~\ref{ch:method} & 
    PATHCAST method \\
  & & Ch.~\ref{ch:implementation} & 
    Python tool implementation \\
\midrule
4 & Demonstration & Ch.~\ref{ch:evaluation} & 
    Application to 48 repositories \\
\midrule
5 & Evaluation & Ch.~\ref{ch:evaluation} & 
    E1: Pipeline Validation \\
  & & & E2: Forecasting Accuracy \\
\midrule
6 & Communication & This thesis & 
    Doctoral dissertation \\
  & & Venues: EASE, MSR & 
    Conference/journal papers \\
\bottomrule
\end{tabular}
\end{table}


\subsection{DSR Knowledge Contribution}
\label{sec:dsr-contribution}

The contribution of this research lies in the formal integration of established techniques into a unified process-oriented prediction pipeline, tailored to the SDLC domain and supported by both theoretical guarantees and empirical validation.

Following the DSR knowledge contribution framework of Gregor and 
Hevner~\cite{gregor2013positioning}, this research operates in the 
\emph{Improvement} quadrant: it addresses a known problem class (software 
process forecasting) with a new solution (process-aware stochastic 
simulation parameterized by process mining). The contribution is not the
invention of new individual techniques (process mining, Markov chains, 
Monte Carlo simulation, and machine learning are established methods), 
but rather their formal integration into a pipeline specifically designed 
for the SDLC domain, accompanied by empirical evidence of its 
effectiveness.


% =============================================================================
% 5.2 DATASET
% =============================================================================
\section{Dataset Description}
\label{sec:dataset}

The empirical evaluation uses a curated dataset of 48 open-source software 
repositories. This section describes the selection criteria, data 
collection process, and descriptive statistics.


\subsection{Repository Selection Criteria}
\label{sec:selection-criteria}

Repositories were selected according to the following inclusion criteria, 
designed to ensure sufficient data volume and process maturity for 
meaningful process mining analysis:

\begin{enumerate}
    \item \textbf{Minimum commit volume:} At least 1{,}000 commits in the 
    main branch, ensuring sufficient event data for Markov estimation 
    with adequate statistical power.
    
    \item \textbf{Issue tracker usage:} Active use of an issue tracking 
    system (GitHub Issues or Jira) with at least 100 closed issues 
    containing state transition data, enabling SDLC process reconstruction.
    
    \item \textbf{Pull request workflow:} Evidence of a pull request-based 
    development workflow, ensuring that code review events are available 
    for event log construction.
    
    \item \textbf{Activity period:} At least 12 months of continuous 
    development activity, providing sufficient temporal depth for 
    temporal hold-out validation.
    
    \item \textbf{Non-trivial contributor base:} At least 5 distinct 
    contributors, excluding single-developer projects where process 
    dynamics are degenerate.
    
    \item \textbf{Non-fork:} The repository is not a fork of another 
    project, avoiding duplicate process data.
\end{enumerate}

Exclusion criteria remove repositories where process mining would produce 
unreliable results:

\begin{enumerate}
    \item Repositories with predominantly bot-generated commits 
    (>50\% of total commits from automated accounts).
    
    \item Repositories where issue tracking is used only for external 
    bug reports (no internal workflow state transitions).
    
    \item Mirror or archive repositories with no active development.
\end{enumerate}


\subsection{Selection Process}
\label{sec:selection-process}

The selection followed a four-step protocol designed to be reproducible
and auditable, with each step documented in the replication
package~\cite{oliveira2025replication}.

\paragraph{Step 1 – Candidate generation.}
An initial candidate list of approximately 150 repositories was
generated using the SEART GitHub Search engine~\cite{seart2021} with
the following query parameters: at least 1,000 commits, issues enabled,
pull requests present, fork status false, last push after 2022-01-01,
and primary language in \{Python, JavaScript/TypeScript, Java, Go, Rust,
C/C++\}. SEART was chosen because it exposes filter parameters not
available through the native GitHub search API and produces a
reproducible, timestamped export~\cite{dabic2021sampling}.

\paragraph{Step 2 – Criterion verification.}
Each candidate was verified against criteria C2 and C5 through direct
calls to the GitHub REST API~\cite{github2023api}. For C2, the issues
timeline events endpoint was queried to confirm the presence of at least
100 closed issues with recorded state transitions. For C5, the
contributors endpoint was used to count distinct contributors, excluding
accounts identified as bots by the GitHub platform. Repositories
failing either criterion were removed from the candidate list.

\paragraph{Step 3 – Stratified selection.}
From the verified candidates, 48 repositories were selected by
stratified sampling to ensure coverage of all dimensions defined in
Table~\ref{tab:dataset-diversity}. The stratification targets were:
three to four repositories per primary language (six language strata);
balanced representation across scale tiers (Small, Medium, Large); and
at least four repositories per governance model (Foundation-backed,
corporate-backed, community-driven). Where multiple candidates competed
for the same stratification cell, preference was given to the repository
with the longer activity span and higher contributor diversity, used
here as proxies for process maturity.

\paragraph{Step 4 – Jira enrichment for Apache and Spring projects.}
For repositories associated with Apache Software Foundation or Spring
Framework projects, the issue tracking component was enriched with data
from the Public Jira Dataset of Montgomery et
al.~\cite{montgomery2022jira}, which provides 32 million historical
state-change records for 1,822 projects across 16 public Jira instances.
The Jira and GitHub issue records were linked by project key, and the
richer Jira transition history was used in place of the GitHub Issues
timeline where available. This substitution is documented
per-repository in Appendix~\ref{app:repository-list} and declared as a
data source variant in the event log metadata field
\texttt{issue\_source}.

The complete list of 48 repositories, with metadata on language, domain,
scale, governance model, workflow maturity, and issue source, is
provided in Appendix~\ref{app:repository-list}.


\subsection{Data Collection}
\label{sec:data-collection}

Data collection is scheduled for completion prior to the empirical
evaluation reported in Chapter~\ref{ch:evaluation}; timestamps will be
updated upon completion. The collection protocol was executed through the
GitHub REST and GraphQL APIs~\cite{github2023api}, with Jira data
sourced from Montgomery et al.~\cite{montgomery2022jira} for the subset
of projects described in Step~4 of
Section~\ref{sec:selection-process}. The following artifacts were
extracted per repository:

\begin{itemize}
    \item \textbf{Git commit history.} Full commit log including hash,
    author, committer, timestamps, message, and file change statistics
    (additions, deletions, files changed), retrieved via the
    \texttt{/repos/\{owner\}/\{repo\}/commits} endpoint with pagination.

    \item \textbf{Issue data.} Complete issue lifecycle including
    creation, state transitions via the timeline events API
    (\texttt{/repos/\{owner\}/\{repo\}/issues/\{number\}/timeline}),
    labels, assignees, milestones, and closure events. For Apache and
    Spring projects, this component was replaced by the Jira transition
    history as described above.

    \item \textbf{Pull request data.} PR creation, review requests,
    review decisions, status check runs, merge events, and linked issue
    references, retrieved via the GraphQL API to minimise request volume.

    \item \textbf{CI/CD data.} GitHub Actions workflow runs including
    trigger event type, status, conclusion, and duration, retrieved via
    the \texttt{/repos/\{owner\}/\{repo\}/actions/runs} endpoint. Only
    repositories with at least 50 recorded workflow runs were classified
    as \textit{Advanced} workflow maturity in
    Table~\ref{tab:dataset-diversity}.
\end{itemize}

All data was collected in JSON format and stored in a normalised
relational schema (PostgreSQL) for efficient querying during event log
construction in Stage~1 of the PATHCAST pipeline.


\subsection{Dataset Descriptive Statistics}
\label{sec:dataset-stats}

The dataset construction strategy balances statistical power, process observability, and diversity, ensuring both internal validity for stochastic modeling and external validity for generalization across different SDLC contexts.

Table~\ref{tab:dataset-summary} presents aggregate statistics for the 
48-repository dataset.

\begin{table}[htbp]
\centering
\caption{Aggregate dataset statistics (48 repositories).}
\label{tab:dataset-summary}
\begin{tabular}{lrrrr}
\toprule
\textbf{Metric} & \textbf{Total} & \textbf{Mean} & \textbf{Median} & \textbf{Std Dev} \\
\midrule
Commits & 227{,}000 & 4{,}729 & 3{,}215 & -- \\
Closed issues & -- & -- & -- & -- \\
Pull requests & -- & -- & -- & -- \\
Contributors & -- & -- & -- & -- \\
Activity span (months) & -- & -- & -- & -- \\
\bottomrule
\multicolumn{5}{l}{\parbox{13cm}{\footnotesize Note: Cells marked ``--'' are
    populated upon completion of data collection
    (Chapter~\ref{ch:evaluation}). Of the 48 planned repositories,
    45 have been confirmed at the time of writing; three slots remain
    pending stratified selection (Step~3,
    Section~\ref{sec:selection-process}). The 227K commits figure
    is confirmed from the SEART candidate-generation phase.}}
\end{tabular}
\end{table}

Table~\ref{tab:dataset-diversity} characterizes the diversity of the 
dataset across dimensions relevant to generalizability.

\begin{table}[htbp]
\centering
\caption{Dataset diversity profile.}
\label{tab:dataset-diversity}
\begin{tabular}{ll}
\toprule
\textbf{Dimension} & \textbf{Coverage} \\
\midrule
Primary languages & Python, JavaScript/TypeScript, Java, Go, Rust, C/C++ \\
Domain categories & Web frameworks, CLI tools, libraries, DevOps, data processing \\
Project scale & Small ($<$5K commits), Medium (5K--20K), Large ($>$20K) \\
Governance model & Foundation-backed, corporate-backed, community-driven \\
Workflow maturity & Basic (issues only), Standard (issues + PRs), Advanced (+ CI/CD) \\
Issue source       & \parbox[t]{7.5cm}{GitHub Issues only, Jira
                     (Apache/Spring enrichment via
                     \cite{montgomery2022jira})} \\
\bottomrule
\end{tabular}
\end{table}

Of the 48 planned repositories, 45 have been confirmed at the time of
writing. The three remaining slots will be filled upon completion of the
stratified selection described in Step~3 of
Section~\ref{sec:selection-process}, prioritising languages or governance
models with fewer than three current representatives. All statistical
analyses reported in Chapter~\ref{ch:evaluation} are based on the final
48-repository dataset; any deviation from this target will be declared
as a limitation in Section~\ref{sec:threats-external}.

The dataset is intentionally heterogeneous to support generalisability
claims. The diversity across languages, domains, scales, and governance
models enables the evaluation to assess whether the proposed method 
generalizes beyond specific project types.


\subsection{Ethical Considerations}
\label{sec:ethics}

All data used in this research is publicly available on GitHub under 
open-source licenses. No private repositories, proprietary data, or 
personally identifiable information beyond publicly visible GitHub 
usernames are used. The research complies with the GitHub Terms of Service 
for API usage and does not involve human subjects.


% =============================================================================
% 5.3 EXPERIMENTAL DESIGN
% =============================================================================
\section{Experimental Design}
\label{sec:eval-design}

The evaluation is structured into two complementary stages, each 
addressing specific research questions. Figure~\ref{fig:eval-design} 
illustrates the experimental design.

\begin{figure}[htbp]
\centering
\begin{tikzpicture}[
    node distance=1.5cm,
    stage/.style={rectangle, draw, thick, rounded corners=4pt, 
                  minimum width=5cm, minimum height=1.4cm, 
                  align=center, font=\small},
    rq/.style={rectangle, draw, dashed, rounded corners=2pt,
               minimum width=2.2cm, minimum height=0.8cm,
               align=center, font=\scriptsize, fill=gray!5},
    arrow/.style={->, thick, >=stealth},
    brace/.style={decorate, decoration={brace, amplitude=8pt, mirror}}
]
    % E1
    \node[stage, fill=blue!10] (e1) {
        \textbf{E1: Pipeline Validation}\\
        {\scriptsize Event log quality + Process model quality}
    };
    \node[rq, right=1cm of e1] (rq1) {RQ1};
    \draw[arrow] (e1) -- (rq1);
    
    % E2
    \node[stage, fill=orange!10, below=1.5cm of e1] (e2) {
        \textbf{E2: Forecasting Accuracy}\\
        {\scriptsize Stochastic forecasts + ML comparison}
    };
    \node[rq, right=1cm of e2, yshift=0.4cm] (rq2) {RQ2};
    \node[rq, right=1cm of e2, yshift=-0.4cm] (rq3) {RQ3};
    \draw[arrow] (e2) -- (rq2);
    \draw[arrow] (e2) -- (rq3);
    
    % Flow
    \draw[arrow] (e1) -- (e2) node[midway, right, font=\scriptsize] 
        {Validated event logs};
    
    % Dataset
    \node[rectangle, draw, fill=green!8, rounded corners=2pt, 
          minimum width=3cm, font=\scriptsize, align=center,
          left=1.5cm of e1, yshift=-0.75cm] (data) {
        48 Repositories\\227K commits\\Dataset $\mathcal{R}$
    };
    \draw[arrow] (data) |- (e1);
    \draw[arrow] (data) |- (e2);
    
\end{tikzpicture}
\caption{Experimental design: two evaluation stages addressing three 
research questions.}
\label{fig:eval-design}
\end{figure}


\subsection{E1: Pipeline Validation}
\label{sec:e1-design}

E1 ensures that subsequent stochastic modelling is not performed on
low-quality process representations, thereby mitigating error propagation
across pipeline stages.

The first evaluation stage assesses whether the PATHCAST pipeline 
(Stages 1 and 2) produces event logs and process models of sufficient 
quality to support stochastic modeling. This stage addresses \textbf{RQ1}.

\subsubsection{Scope}

E1 executes Stages 1 and 2 of the pipeline on all 48 repositories 
and evaluates the outputs:

\begin{enumerate}
    \item \textbf{Event log quality:} Measured by correlation coverage 
    (fraction of raw events successfully assigned to cases), activity 
    coverage (fraction of the SDLC taxonomy represented in the log), 
    and trace completeness (fraction of cases reaching an absorbing 
    state).
    
    \item \textbf{Process model quality:} Measured by fitness, precision, 
    and generalization of the discovered models, using the alignment-based 
    conformance checking described in Section~\ref{sec:s2-conformance}.
\end{enumerate}

\subsubsection{Protocol}

For each repository $r_i \in \mathcal{R}$:

\begin{enumerate}
    \item Execute Stage 1 (Event Log Construction) with the default 
    correlation configuration.
    
    \item Compute event log quality metrics 
    (Definition~\ref{def:log-quality}).
    
    \item Execute Stage 2 (Process Mining) with the Inductive Miner 
    Infrequent variant at noise threshold $\theta_{\text{noise}} = 0.2$.
    
    \item Compute process model quality metrics (fitness, precision, 
    generalization).
    
    \item Record the number of unique process variants and the state
    space $|S|$ for downstream Markov estimation.

    \item Flag state-transition pairs with fewer than 30 observed
    transitions. These pairs receive Laplace additive smoothing
    ($\alpha = 1$) during Markov estimation
    (Section~\ref{sec:s3-mle}), ensuring that the transition
    matrix $\hat{P}$ remains fully defined even for infrequently
    observed transitions. The count of smoothed pairs is reported
    per repository as a data-quality indicator.
\end{enumerate}

\subsubsection{Sensitivity Analysis}

The noise threshold $\theta_{\text{noise}}$ is varied over the range 
$\{0.0, 0.1, 0.2, 0.3, 0.5\}$ on a representative subset of 10 
repositories to assess the sensitivity of discovery quality to this 
parameter. The analysis reports the Pareto front of fitness versus 
precision, identifying the threshold that best balances both objectives.

\subsubsection{Acceptance Criteria}

A repository is considered suitable for stochastic modeling (Stage 3) 
if its event log and process model satisfy:

\begin{equation}
\label{eq:acceptance}
\text{fitness}(\mathcal{L}, M) \geq 0.7 \;\wedge\;
\text{coverage}(\mathcal{L}) \geq 0.5 \;\wedge\;
|S_T| \geq 2 \;\wedge\;
|\mathcal{L}_{\text{test}}| \geq 50
\end{equation}

The fourth condition, $|\mathcal{L}_{\text{test}}| \geq 50$, requires
that the temporal hold-out test set
(Definition~\ref{def:temporal-holdout}) contains at least 50 completed
cases. Below this threshold, empirical percentile estimates exhibit
high variance and CRPS becomes unreliable~\cite{gneiting2007strictly}.
Repositories that satisfy the first three conditions but fail the case
count threshold are reported separately as \emph{process-quality
validated but statistically insufficient for forecasting comparison};
their E1 results are retained in full while they are excluded from E2
only. Repositories failing any of the four criteria are excluded from
E2 and reported separately, with analysis of the corresponding failure
modes.


\subsection{E2: Forecasting Accuracy}
\label{sec:e2-design}

The second evaluation stage assesses the accuracy and calibration of 
the probabilistic forecasts produced by Stages 3 and 4, and evaluates 
the predictive contribution of process-derived ML features. This stage 
addresses \textbf{RQ2} and \textbf{RQ3}.

\subsubsection{Scope}

E2 operates on the repositories that passed all four acceptance
criteria of E1, including the minimum test-set size of 50~completed
cases (Equation~\ref{eq:acceptance}). Repositories excluded from E2
for failing the case count threshold are listed in the replication
package with their E1 quality metrics, ensuring full transparency
about the effective sample size for each RQ. For each qualifying
repository:

\begin{enumerate}
    \item Execute Stages 3 and 4 (Markov Estimation + Monte Carlo
    Simulation) to produce probabilistic lead time forecasts.

    \item Execute the ML component with three feature sets (PM-only,
    MSR-only, Combined) to produce case-level outcome predictions.

    \item Compare all forecasts against ground truth (observed lead
    times and outcomes).
\end{enumerate}

\subsubsection{Temporal Hold-Out Protocol}
\label{sec:temporal-holdout}

To simulate realistic deployment conditions, where a model is trained 
on historical data and used to forecast future outcomes, a temporal 
hold-out split is applied:

\begin{definition}[Temporal Hold-Out]
\label{def:temporal-holdout}
For each repository with event log $\mathcal{L}$ spanning the time 
interval $[t_{\min}, t_{\max}]$:
\begin{align}
t_{\text{cut}} &= t_{\min} + 0.7 \cdot (t_{\max} - t_{\min}) 
    \label{eq:cutoff} \\
\mathcal{L}_{\text{train}} &= \{c \in \mathcal{L} : 
    \max(\tau(\sigma_c)) \leq t_{\text{cut}}\} 
    \label{eq:train-set} \\
\mathcal{L}_{\text{test}} &= \{c \in \mathcal{L} : 
    \min(\tau(\sigma_c)) > t_{\text{cut}}\} 
    \label{eq:test-set}
\end{align}

The training set contains cases completed before the cutoff date. The 
test set contains cases that started after the cutoff date, ensuring 
no temporal overlap.
\end{definition}

The 70/30 split ratio is chosen as a standard convention in temporal 
validation for software engineering prediction 
tasks~\cite{tashman2000out}. Sensitivity to the split ratio is assessed 
by repeating the evaluation with 60/40 and 80/20 splits on a 
representative subset.

\textbf{Model estimation:} The Markov transition matrix $P$, sojourn 
distributions $\hat{F}$, and ML models are estimated exclusively from 
$\mathcal{L}_{\text{train}}$.

\textbf{Forecast evaluation:} Forecasts are generated for each case 
$c \in \mathcal{L}_{\text{test}}$ and compared against the observed 
lead time $T_c^{\text{obs}}$ and outcome $y_c^{\text{obs}}$.


\subsubsection{Cross-Repository Generalization}
\label{sec:cross-repo}

To assess whether the method generalizes across projects, a
leave-one-project-out (LOPO) evaluation is conducted on subsets of
repositories within the same domain (e.g., all Python web frameworks):

\begin{enumerate}
    \item Train the Markov model and ML predictor on $k-1$ repositories.
    \item Generate forecasts for the held-out repository.
    \item Repeat for each repository in the group.
\end{enumerate}

\noindent\textbf{Minimum group size.} The LOPO protocol requires at
least five repositories per domain group to produce a meaningful
leave-one-out sample. Groups with fewer than five repositories are
excluded from LOPO inference and reported with descriptive statistics
only. With the current dataset composition, Python web frameworks
(Django, Flask, FastAPI, Requests, Click~=~5) and Go DevOps tools
(kubectl, Helm, Terraform, Prometheus, cli~=~5) satisfy this
requirement; other domain groups are reported descriptively.

This protocol tests the transferability of transition patterns across
projects with similar development cultures.


%============================================================
% Section 5.4  Evaluation Metrics
% ============================================================

\section{Evaluation Metrics}
\label{sec:evaluation-metrics}
 
The experimental design ensures internal validity through temporal isolation,
external validity through cross-project evaluation, and construct validity
through multiple metrics.
This section defines the metrics used to evaluate the proposed method,
organized by research question.
 
% ------------------------------------------------------------
\subsection{Metrics for RQ1: Event Log and Process Model Quality}
\label{sec:metrics-rq1}
 
Answering RQ1 requires evaluating two distinct aspects of pipeline
output: the quality of the constructed event log and the quality
of the process model discovered from it.
These two aspects are treated separately because degradation in
event log quality propagates to the model and can obscure whether
discovery failures stem from sparse data or from modeling
limitations.
Table~\ref{tab:metrics-rq1} summarizes the full metric set;
the following paragraphs justify each choice.
 
\begin{table}[ht]
\centering
\caption{Evaluation metrics for RQ1.}
\label{tab:metrics-rq1}
\small
\begin{tabular}{@{}l l p{7.8cm}@{}}
\toprule
\textbf{Metric} & \textbf{Range} & \textbf{Definition} \\
\midrule
\multicolumn{3}{@{}l}{\textit{Event Log Quality}} \\[2pt]
Correlation Coverage & $[0, 1]$ &
  Fraction of raw events assigned to a case (Equation~4.15). \\[3pt]
Activity Coverage    & $[0, 1]$ &
  Fraction of SDLC taxonomy activities observed in the log
  (Equation~4.22). \\[3pt]
Trace Completeness   & $[0, 1]$ &
  Fraction of cases reaching an absorbing state (Equation~4.20). \\[3pt]
Mean Trace Length    & $\mathbb{R}^+$ &
  Average number of events per trace (Equation~4.21). \\[3pt]
Variant Count        & $\mathbb{N}$ &
  Number of unique activity sequences (Equation~4.23). \\[4pt]
\midrule
\multicolumn{3}{@{}l}{\textit{Process Model Quality}} \\[2pt]
Fitness              & $[0, 1]$ &
  Alignment-based fitness (Equation~2.2). Higher is better. \\[3pt]
Precision            & $[0, 1]$ &
  Fraction of model-allowed behaviour observed in the log.
  Higher is better. \\[3pt]
Generalization       & $[0, 1]$ &
  Degree to which the model captures the underlying process rather
  than overfitting to the observed traces. Higher is better. \\
\bottomrule
\end{tabular}
\end{table}
 
\paragraph{Correlation Coverage.}
Stage~1 of PATHCAST correlates raw events from heterogeneous sources
(Git commits, issue tracker transitions, CI/CD runs) into unified
case traces.
Correlation Coverage measures the fraction of raw events that were
successfully attributed to a case identifier.
A low value indicates that the cascade correlation strategy failed to
resolve a substantial portion of the raw data, either due to missing
cross-references, incomplete issue keys in commit messages, or gaps in
the CI/CD linkage.
The metric provides a direct signal of Stage~1 output fidelity and is
the first indicator of whether the pipeline can produce a usable log
from a given repository.
 
\paragraph{Activity Coverage.}
Not all repositories exercise every SDLC phase defined in the reference
taxonomy (Definition~\ref{def:activity-taxonomy}).
Activity Coverage measures the proportion of taxonomy-defined activities
that appear at least once in the constructed log.
A repository that exercises few taxonomy activities will yield a
Markov model with a sparse state space, reducing the generality of
downstream forecasts.
The metric thus contextualizes fitness and precision results: a high
fitness on a log with low activity coverage may reflect an overly
simple model rather than genuine predictive power.
 
\paragraph{Trace Completeness.}
A trace is complete if it reaches one of the designated absorbing
states (Done or Cancelled).
Incomplete traces correspond to work items that are still in progress
at the time of data collection, or whose final state could not be
determined from the available sources.
Trace Completeness directly bounds the size of the training set
available for Markov estimation and ML training, since only complete
traces contribute to the supervised learning protocol
(Section~\ref{sec:temporal-holdout}).
Repositories with low Trace Completeness are flagged and excluded
from the comparison analysis if the count of complete cases falls
below the minimum threshold $n_{\min}^{\text{cases}} = 50$.
 
\paragraph{Mean Trace Length and Variant Count.}
These two metrics characterize the structural complexity of the
discovered process.
Mean Trace Length indicates the average number of observed state
transitions per case and serves as a proxy for process depth: shallow
processes generate short traces that provide limited evidence for
Markov estimation, whereas deep processes may indicate fragmented
correlation strategies that split single work items into many events.
Variant Count enumerates the distinct activity sequences observed
across all traces.
A high variant count relative to the number of traces signals process
diversity or noise in activity labeling, both of which affect the
reliability of the discovered process model.
Together, these two metrics inform the interpretation of fitness and
precision results reported in Evaluation Stage~E1.
 
\paragraph{Process Model Quality: Fitness, Precision, and Generalization.}
Once the event log passes the quality threshold, the process model
discovered by the Inductive Miner is evaluated on three
complementary dimensions.
Fitness, computed via alignment-based conformance checking
(Equation~2.2), measures the extent to which the model
can replay the observed traces without deviations.
A fitness value close to~1 indicates that the model accounts for the
full behavioral repertoire present in the log.
Precision penalizes the opposite failure mode: a model that allows
many behaviors not present in the log is imprecise, even if it
perfectly replays all observed traces.
Overly imprecise models produce permissive Markov transition matrices
that inflate forecast uncertainty by assigning probability mass to
transitions that never occur in practice.
Generalization assesses whether the model captures the underlying
process rather than overfitting the specific traces in the log.
A model with high fitness and high precision but low generalization
has memorized the training traces; it will perform poorly on new cases
that follow legitimate but previously unobserved paths.
 
The combination of fitness, precision, and generalization provides a
balanced assessment consistent with the four-quality-dimension
framework of van der Aalst~\cite{aalst2016process}, which
additionally includes simplicity as a fourth dimension.
Simplicity is not reported as a primary metric in this evaluation
because the state space is constrained by the SDLC taxonomy
(Section~\ref{sec:s1-mapping}), limiting model complexity by
construction.
 
% ------------------------------------------------------------
\subsection{Metrics for RQ2: Stochastic Forecasting Accuracy}
\label{sec:metrics-rq2}
 
Evaluating probabilistic forecasts requires metrics that assess both
point accuracy and distributional calibration.
Table~\ref{tab:metrics-rq2} lists the full metric set for RQ2.
 
\begin{table}[ht]
\centering
\caption{Evaluation metrics for RQ2.}
\label{tab:metrics-rq2}
\small
\begin{tabular}{@{}l l p{7.2cm}@{}}
\toprule
\textbf{Metric} & \textbf{Type} & \textbf{Definition} \\
\midrule
\multicolumn{3}{@{}l}{\textit{Point Forecast Accuracy}} \\[2pt]
MAE  & Error &
  $\dfrac{1}{n}\displaystyle\sum_c
    \lvert T_c^{\text{obs}} - \hat{T}_c^{\text{P50}} \rvert$ \\[8pt]
MAPE & Error (\%) &
  $\dfrac{100}{n}\displaystyle\sum_c
    \dfrac{\lvert T_c^{\text{obs}} - \hat{T}_c^{\text{P50}} \rvert}
          {T_c^{\text{obs}}}$ \\[8pt]
MdAE & Error &
  $\operatorname{median}
    \!\left(\lvert T_c^{\text{obs}} - \hat{T}_c^{\text{P50}} \rvert\right)$ \\[6pt]
\midrule
\multicolumn{3}{@{}l}{\textit{Distributional Accuracy}} \\[2pt]
CRPS & Score &
  $\dfrac{1}{n}\displaystyle\sum_c
    \int_{-\infty}^{\infty}
    \!\left[\hat{F}_c(t) - \mathbf{1}(t \geq T_c^{\text{obs}})\right]^2
    \mathrm{d}t$ \\[10pt]
Calibration Error & Error &
  $\mathrm{ACE} =
    \dfrac{1}{9}\displaystyle\sum_{\alpha \in \{0.1,\ldots,0.9\}}
    \lvert \hat{p}_\alpha - \alpha \rvert$ \\[8pt]
Coverage (85\%)   & $[0,1]$ &
  Fraction of $T_c^{\text{obs}}$ inside
  $[\hat{Q}_{0.075},\,\hat{Q}_{0.925}]$. \\[4pt]
\midrule
\multicolumn{3}{@{}l}{\textit{Outcome Prediction}} \\[2pt]
Completion Accuracy & $[0,1]$ &
  Accuracy of predicting Done vs.\ Cancelled outcome. \\
\bottomrule
\end{tabular}
\end{table}
 
\paragraph{CRPS.}
The Continuous Ranked Probability Score (CRPS) is the primary metric
for RQ2.
Unlike MAE, which evaluates only the point forecast (P50), CRPS
evaluates the entire predicted distribution, penalizing both
inaccuracy and overconfidence.
A lower CRPS indicates a better probabilistic forecast.
CRPS reduces to MAE when the forecast degenerates to a point estimate,
preserving comparability with deterministic baselines
\cite{gneiting2007strictly}.
This property makes CRPS the preferred metric whenever both point and
distributional forecasts are compared in the same study.
 
\paragraph{Point Forecast Metrics: MAE, MAPE, and MdAE.}
Although CRPS subsumes MAE as a special case, reporting MAE and MAPE
separately enables direct comparison with prior software effort
estimation studies that use only point forecasts.
MAE is reported in calendar days and represents the average absolute
deviation between the P50 forecast and the observed lead time.
MAPE expresses the same deviation as a percentage of the observed
value, facilitating comparisons across repositories with different
characteristic lead-time scales.
MdAE complements MAE by providing a robust central tendency estimate
that is less sensitive to outliers, which are common in software
lead-time distributions due to blocked or deprioritized work items.
 
\paragraph{Calibration Error.}
Calibration Error assesses whether the predicted percentiles match
observed frequencies.
A well-calibrated forecast at the 85th percentile should contain
approximately 85\% of observed values below that threshold.
Aggregate Calibration Error (ACE) averages the absolute deviation
across nine decile checkpoints, from $\alpha = 0.1$ to
$\alpha = 0.9$:
\begin{equation}
  \mathrm{ACE} = \frac{1}{9}
    \sum_{\alpha \in \{0.1, 0.2, \ldots, 0.9\}}
    \lvert \hat{p}_\alpha - \alpha \rvert
\end{equation}
An ACE of zero denotes perfect calibration; values above 0.1 indicate
systematic over- or underconfidence.
ACE is reported alongside CRPS in all experiments because a forecast
can achieve a low CRPS while remaining miscalibrated if the
distributional shape is correct but the percentile coverage is
skewed.
 
\paragraph{Coverage at 85\%.}
The 85th-percentile interval $[\hat{Q}_{0.075},\,\hat{Q}_{0.925}]$
is the standard agile commitment interval recommended by Vacanti
\cite{vacanti2015}.
Coverage at 85\% is a single-interval calibration check that
practitioners can relate directly to delivery commitments: if the
method claims that 85\% of work items complete within a given
window, empirical coverage should approximate that claim.
This metric complements ACE by focusing on the single interval most
relevant to operational planning, rather than aggregating across
all deciles.
 
\paragraph{Completion Accuracy.}
In addition to lead-time forecasts, PATHCAST predicts whether a work
item will reach Done or Cancelled (Definition~\ref{def:label}).
Completion Accuracy measures the proportion of cases for which the
predicted terminal outcome matches the observed outcome.
This metric is relevant when decision-makers need to estimate scope at
risk, that is, the fraction of planned work items that will not
complete successfully.
 
% ------------------------------------------------------------
\subsection{Metrics for RQ3: ML Feature Comparison}
\label{sec:metrics-rq3}
 
RQ3 asks whether process-derived features add predictive value
beyond traditional repository metrics.
Answering this question requires classification metrics that
assess both discrimination and calibration, as well as an
interpretable feature-ranking mechanism.
Table~\ref{tab:metrics-rq3} defines the metric set.
 
\begin{table}[ht]
\centering
\caption{Evaluation metrics for RQ3.}
\label{tab:metrics-rq3}
\small
\begin{tabular}{@{}l l p{7.4cm}@{}}
\toprule
\textbf{Metric} & \textbf{Type} & \textbf{Definition} \\
\midrule
AUC-ROC &
  Discrimination &
  Area under the Receiver Operating Characteristic (ROC) curve.
  Threshold-independent discrimination between delayed and
  on-time outcomes. \\[4pt]
F1-Score &
  Balance &
  $2 \cdot \dfrac{\text{Precision} \cdot \text{Recall}}
                  {\text{Precision} + \text{Recall}}$
  Computed at the threshold maximising the harmonic mean on the
  validation set. \\[8pt]
Brier Score &
  Calibration &
  $\dfrac{1}{n}\displaystyle\sum_c (\hat{p}_c - y_c)^2$.
  Lower is better. Assesses probability calibration quality. \\[6pt]
Feature Importance &
  Ranking &
  Permutation importance (primary) and Gini importance from
  tree-based models (RF, XGBoost). \\
\bottomrule
\end{tabular}
\end{table}
 
\paragraph{AUC-ROC.}
The Area Under the Receiver Operating Characteristic Curve (AUC-ROC)
is the primary discrimination metric for the binary outcome predictor
(on-time vs.\ delayed, as defined in
Definition~\ref{def:label}).
AUC-ROC is threshold-independent: it evaluates the model's ability
to rank positive instances above negative ones across all possible
classification thresholds, making it robust to the class imbalance
that commonly arises when a small fraction of work items is
categorized as delayed under a given threshold strategy.
A value of 0.5 corresponds to random classification; a value of 1.0
denotes perfect discrimination.
 
\paragraph{F1-Score.}
F1-Score provides a threshold-dependent complement to AUC-ROC.
Because the cost of missing a delayed item typically exceeds the
cost of a false alarm in practice, the F1-Score is computed at the
threshold that maximizes the harmonic mean of precision and recall
on the validation set.
This threshold is then applied to the test set without adjustment,
preserving the temporal integrity of the evaluation.
 
\paragraph{Brier Score.}
The Brier Score evaluates the calibration of the probabilistic
classifier, measuring the mean squared deviation between the
predicted delay probability $\hat{p}_c$ and the binary outcome
label $y_c$.
A Brier Score of zero indicates perfect calibration; the Brier Score
for a random classifier equals the outcome class frequency times one
minus that frequency.
Reporting the Brier Score alongside AUC-ROC distinguishes models
that correctly rank cases from models that additionally produce
reliable probability estimates, a distinction relevant to risk-based
portfolio management.
 
\paragraph{Feature Importance.}
Feature importance rankings determine the relative contribution of
each feature group (process-derived vs.\ traditional MSR) to
predictive performance.
For tree-based models (Random Forest and XGBoost), both Gini
importance and permutation importance are computed; permutation
importance is reported as the primary ranking because it is
invariant to the scale and cardinality of individual features.
For Logistic Regression, coefficients are standardized before
reporting.
 
The three feature sets defined in Section~\ref{sec:ml-feature-eng}
are compared:
\begin{enumerate}
  \item \textbf{PM-only}: Features~1--8 (process-derived).
  \item \textbf{MSR-only}: Features~9--15 (traditional repository metrics).
  \item \textbf{Combined}: Features~1--15 (all features).
\end{enumerate}
A statistically meaningful improvement of PM-only over MSR-only, or
of Combined over MSR-only, would demonstrate that process
mining-derived features carry predictive information not captured by
traditional metrics, answering RQ3 affirmatively.
The Wilcoxon signed-rank test is applied across the 48 repositories
to assess statistical significance; effect size is quantified with
Cliff's~$\delta$ following Kampenes et al.\
\cite{kampenes2007systematic}.
% =============================================================================
% 5.5 BASELINES
% =============================================================================
\section{Baseline Methods}
\label{sec:baselines}

The proposed process-aware Monte Carlo method is compared against four baselines, each representing a distinct level of process awareness, and this baseline hierarchy enables controlled ablation of the modeling components, isolating the contribution of process awareness, stochastic simulation, and distributional forecasting.

\begin{table}[htbp]
\centering
\caption{Baseline methods for comparison.}
\label{tab:baselines}
\small
\begin{tabular}{lp{4cm}p{5cm}}
\toprule
\textbf{Baseline} & \textbf{Description} & \textbf{Rationale} \\
\midrule
B1: Historical Mean & 
    Point forecast: mean lead time of completed cases in 
    $\mathcal{L}_{\text{train}}$. & 
    Na\"ive baseline. Any reasonable method should outperform this. \\
\midrule
B2: Throughput MC & 
    Throughput-based Monte Carlo 
    simulation~\cite{vacanti2015, savage2009flaw,magennis2011forecasting}. Samples daily 
    throughput from historical data. & 
    State-of-practice in agile forecasting. Black-box: no 
    process structure. \\
\midrule
B3: Markov Analytical & 
    Expected lead time from the fundamental matrix: 
    $\mathbb{E}[\text{LT}] = \sum_j N_{ij} \cdot \bar{d}_j$. & 
    Process-aware but deterministic (point estimate only). 
    Tests the marginal value of Monte Carlo simulation over 
    analytical computation. \\
\midrule
B4: Historical Distribution & 
    Empirical CDF of lead times from 
    $\mathcal{L}_{\text{train}}$. No modeling; direct 
    resampling from observed lead times. & 
    Distributional baseline. Tests whether the stochastic model 
    adds value beyond empirical resampling. \\
\bottomrule
\end{tabular}
\end{table}

The baselines are ordered by increasing sophistication: B1 (no model), 
B4 (empirical distribution), B2 (stochastic, no process structure), 
B3 (process-aware, no simulation), and the proposed method (process-aware 
stochastic simulation). If the proposed method outperforms B2, the 
value of process awareness is demonstrated. If it outperforms B3, the 
value of distributional forecasts over point estimates is demonstrated. 
If it outperforms B4, the value of the stochastic process model over 
na\"ive resampling is demonstrated.


% =============================================================================
% 5.6 STATISTICAL ANALYSIS
% =============================================================================
\section{Statistical Analysis Procedures}
\label{sec:statistical-analysis}

\subsection{Within-Repository Comparison}
\label{sec:within-repo}

For each repository, the proposed method and each baseline produce 
forecasts for the same set of test cases $\mathcal{L}_{\text{test}}$. 
The comparison uses paired statistical tests because the same cases are 
evaluated across methods.

\textbf{Normality assessment:} The Shapiro-Wilk test is applied to the 
distribution of per-case error differences between the proposed method 
and each baseline. If normality holds ($p > 0.05$), a paired $t$-test 
is used. Otherwise, the Wilcoxon signed-rank test is applied.

\textbf{Effect size:} Cohen's $d$ (for parametric comparisons) or 
Cliff's $\delta$ (for non-parametric comparisons) is reported alongside 
$p$-values, following the recommendations of 
Kampenes et al.~\cite{kampenes2007systematic} for empirical software 
engineering studies.

Statistical significance is complemented by the analysis of effect size, allowing one to distinguish the practical relevance of differences that are merely statistical.

\subsection{Cross-Repository Aggregation}
\label{sec:cross-repo-stats}

To aggregate results across the 48 repositories, the following approach 
is adopted:

\begin{enumerate}
    \item For each repository $r_i$, compute the per-repository metric 
    (e.g., CRPS) for the proposed method and each baseline.
    
    \item Treat the 48 per-repository metric values as paired samples 
    (one pair per repository, proposed vs. baseline).
    
    \item Apply the Wilcoxon signed-rank test across repositories to 
    assess whether the proposed method outperforms each baseline 
    consistently.
    
    \item Report the median improvement and interquartile range across 
    repositories.
\end{enumerate}

\textbf{Multiple comparison correction:} Because the proposed method is 
compared against four baselines, the Bonferroni correction is applied, 
adjusting the significance level to $\alpha_{\text{adj}} = 0.05 / 4 = 0.0125$.


\subsection{Statistical Power Analysis}
\label{sec:power-analysis}

The adequacy of the 48-repository dataset for the planned statistical
comparisons is assessed through an a priori power analysis following
the recommendations of Kampenes et al.~\cite{kampenes2007systematic}
for empirical software engineering experiments.

\textbf{RQ2 and RQ3 — Cross-repository comparison (Wilcoxon
signed-rank, $n = 48$).}
The primary statistical test uses 48 paired per-repository observations
(proposed method vs.\ each baseline). Under the Bonferroni-adjusted
significance level $\alpha_{\text{adj}} = 0.05 / 4 = 0.0125$, the
achieved power for a medium effect size ($r = 0.3$) is
$1 - \beta \approx 0.72$, and for a large effect ($r = 0.5$) it
exceeds 0.97. The expected effect sizes are informed by Rogge-Solti
and Weske~\cite{RoggeSolti2015NonMarkovianSPN}, who reported large
differences between stochastic process-aware models and exponential
baselines on comparable BPM datasets. Power is computed via:

\begin{equation}
\label{eq:power}
1 - \beta =
P\!\left(Z > z_{\alpha/2} - r\sqrt{n-1}\right) +
P\!\left(Z < -z_{\alpha/2} - r\sqrt{n-1}\right)
\end{equation}

\noindent where $Z \sim \mathcal{N}(0,1)$, $r$ is the point-biserial
effect size, $n = 48$, and $z_{\alpha/2} = z_{0.00625} \approx 2.50$
(Bonferroni-adjusted). Table~\ref{tab:power-analysis} summarises power
estimates across effect sizes.

\begin{table}[htbp]
\centering
\caption{A priori power estimates for the cross-repository Wilcoxon
  signed-rank test ($n = 48$, four baselines).}
\label{tab:power-analysis}
\small
\begin{tabular}{lccc}
\toprule
\textbf{Effect size} & \textbf{$r$} &
  \textbf{Power ($\alpha = 0.05$)} &
  \textbf{Power ($\alpha_{\text{adj}} = 0.0125$)} \\
\midrule
Small            & 0.1 & 0.17 & 0.09 \\
Small--medium    & 0.2 & 0.44 & 0.27 \\
Medium           & 0.3 & 0.75 & 0.55 \\
Medium--large    & 0.4 & 0.92 & 0.80 \\
Large            & 0.5 & 0.98 & 0.94 \\
\bottomrule
\end{tabular}
\end{table}

The dataset is adequately powered for medium to large effects at the
adjusted significance level. Low power for small effects ($r \leq 0.2$)
is an accepted limitation: the research hypotheses concern practically
meaningful improvements over process-blind baselines, which require
medium or large effect sizes to be scientifically credible.

\textbf{RQ3 — ML within-repository comparison.}
For the repositories satisfying
$|\mathcal{L}_{\text{test}}| \geq 50$, the median expected test-set
size is approximately 80--120 cases, based on the confirmed 227K
commits and a conservative issue-to-commit ratio of~1:4. This is
sufficient for reliable AUC-ROC estimation, which stabilises above
$n \approx 50$ for approximately balanced
datasets~\cite{hanley1982meaning}. Repositories near the threshold are
reported with individual confidence intervals.


\subsection{Reporting Standards}
\label{sec:reporting}
The aggregation strategy ensures robustness to the heterogeneity of the projects, while at the same time preserving statistical power through paired comparisons.

All results are reported following the guidelines of 
Wohlin et al.~\cite{wohlin2012experimentation} for controlled 
experiments in software engineering:

\begin{itemize}
    \item Descriptive statistics: median, IQR, and box plots 
    (preferred over mean and standard deviation for potentially 
    skewed distributions).
    
    \item Statistical tests: test name, test statistic, $p$-value, 
    and effect size for every comparison.
    
    \item Confidence intervals: 95\% confidence intervals for all 
    aggregate metrics.
    
    \item Visualization: calibration plots (predicted vs. observed 
    quantile coverage), CRPS comparison box plots, and feature 
    importance rankings.
\end{itemize}


% =============================================================================
% 5.7 THREATS TO VALIDITY
% =============================================================================
\section{Threats to Validity}
\label{sec:threats}

Following the classification of Wohlin et 
al.~\cite{wohlin2012experimentation}, the threats to validity are 
organized into four categories.


\subsection{Internal Validity}
\label{sec:threats-internal}

\textbf{Temporal information leakage:} The most common threat in 
predictive studies on software repositories is the use of future 
information to predict past outcomes. This is mitigated by the strict 
temporal hold-out protocol (Definition~\ref{def:temporal-holdout}), 
which ensures that training and test sets are temporally disjoint with 
no overlap.

\textbf{Confounding by project phase:} A repository may exhibit 
different process dynamics during early development versus maintenance 
phases. The training set may capture one phase while the test set 
captures another. This is partially mitigated by the 70/30 temporal 
split, which provides a substantial training window, and by the 
sensitivity analysis with alternative split ratios.

\textbf{Case correlation errors:} Incorrect case correlation in Stage 1 
may introduce noise into the event log, degrading downstream model 
quality. This threat is mitigated by the correlation quality metrics 
(Definition~\ref{def:corr-quality}) reported in E1, and by the cascade 
strategy that prioritizes high-confidence explicit references over 
heuristic matching.

\textbf{Bot contamination:} Automated commits from bots (Dependabot, 
Renovate) may distort process models if not filtered. The bot filtering 
step in Stage 1 mitigates this threat, and a sensitivity analysis 
compares results with and without bot filtering on a subset of 
repositories.


\subsection{External Validity}
\label{sec:threats-external}

\textbf{Open-source bias:} The dataset consists exclusively of 
open-source repositories hosted on GitHub. Development practices, 
governance models, and process maturity in open-source projects may 
differ from proprietary industrial projects. The results should be 
interpreted as evidence of feasibility and relative method performance 
rather than as absolute forecasting accuracy applicable to all settings.

\textbf{GitHub platform bias:} All repositories use GitHub as the 
hosting platform. Workflows, PR conventions, and issue tracking 
practices may differ on other platforms (GitLab, Bitbucket, Azure 
DevOps). The event extraction framework (Stage 1) is designed with 
pluggable source extractors to accommodate alternative platforms, 
but this extensibility is not empirically validated in this thesis.

\textbf{Language and domain bias:} Despite deliberate heterogeneity in 
the dataset (Table~\ref{tab:dataset-diversity}), certain languages and 
domains may be overrepresented. The cross-repository generalization 
protocol (Section~\ref{sec:cross-repo}) partially addresses this by 
testing transferability within and across domain groups.


\subsection{Construct Validity}
\label{sec:threats-construct}

\textbf{Lead time definition:} Lead time is measured from the first 
event in a case trace to the last event. This definition may not align 
with organizational definitions of lead time (e.g., from customer 
request to deployment). The choice is justified by data availability 
(the first and last events are observable in all repositories) and 
by consistency with the process mining-based measurement.

\textbf{Delay threshold for ML labels:} The binary classification label 
(on-time vs. delayed) depends on the threshold $\tau_{\text{threshold}}$ 
(Definition~\ref{def:label}). Different thresholds produce different 
class distributions and may affect ML performance. This threat is 
mitigated by evaluating three threshold strategies (percentile-based, 
SLA-based, relative) and reporting results for each.

\textbf{Metric selection.} CRPS is chosen as the primary metric because
it evaluates the entire predicted distribution rather than a single
point estimate, consistent with the thesis hypothesis that distributional
forecasts are more informative than point forecasts. To preserve
comparability with point-forecast baselines, MAE is reported alongside
CRPS in all comparisons.



\subsection{Conclusion Validity}
\label{sec:threats-conclusion}

\textbf{Statistical power:} The number of test cases per repository
varies across the dataset. Repositories near the minimum threshold
($|\mathcal{L}_{\text{test}}| = 50$) produce per-repository CRPS and
calibration estimates with higher variance than larger repositories.
This threat is mitigated by three design choices: (i)~the minimum case
count condition in Equation~\ref{eq:acceptance} excludes repositories
with fewer than 50 test cases from E2; (ii)~the cross-repository
aggregation treats each qualifying repository as a single paired
observation, so the Wilcoxon signed-rank test operates on $n = 48$
pairs, achieving power $1 - \beta \geq 0.72$ for medium effect sizes
at $\alpha_{\text{adj}} = 0.0125$ (Section~\ref{sec:power-analysis},
Table~\ref{tab:power-analysis}); and (iii)~per-repository results are
reported with individual confidence intervals, enabling readers to
assess reliability for each repository independently. The a priori
analysis confirms that effect sizes consistent with those reported in
comparable process-aware forecasting
studies~\cite{RoggeSolti2015NonMarkovianSPN} are detectable with
adequate power.

\textbf{Multiple comparisons:} Comparing the proposed method against 
four baselines on multiple metrics increases the risk of Type I errors. 
The Bonferroni correction (Section~\ref{sec:cross-repo-stats}) 
addresses this by adjusting the significance threshold.

\textbf{Reproducibility:} To support reproducibility, the following 
artifacts are made available: (i) the list of 48 repository identifiers
and collection dates; (ii) the source code of the PATHCAST pipeline;
(iii) the evaluation scripts; and (iv) the intermediate data files
(event logs, transition matrices, simulation results). Raw repository
data is not redistributed but can be independently re-collected via the
GitHub REST API~\cite{github2023api} using the identifiers listed in
Appendix~\ref{app:repository-list}.

Although these threats cannot be completely eliminated, the combination
of design choices, validation protocols, and sensitivity analyses ensures
that their impact is systematically addressed and reported.


% =============================================================================
% 5.8 CHAPTER SUMMARY
% =============================================================================
\section{Chapter Summary}
\label{sec:methodology-summary}

This chapter presented the methodological framework for the empirical 
evaluation of the PATHCAST method. Table~\ref{tab:methodology-summary} 
synthesizes the mapping between research questions, evaluation stages, 
metrics, and baselines.

\begin{table}[htbp]
\centering
\caption{Summary: research questions, evaluation stages, metrics, and baselines.}
\label{tab:methodology-summary}
\small
\begin{tabular}{llll}
\toprule
\textbf{RQ} & \textbf{Stage} & \textbf{Primary Metrics} & \textbf{Baselines} \\
\midrule
RQ1 & E1 & 
    Fitness, precision, coverage & 
    N/A (quality assessment) \\
RQ2 & E2 & 
    CRPS, calibration error, MAE & 
    B1, B2, B3, B4 \\
RQ3 & E2 & 
    AUC-ROC, F1, Brier score & 
    PM-only vs. MSR-only vs. Combined \\
\bottomrule
\end{tabular}
\end{table}

The evaluation is designed to answer three questions with increasing 
specificity: Can the pipeline produce usable event logs and models from 
real repositories? (E1, RQ1). Do process-aware stochastic forecasts 
outperform process-blind baselines? (E2, RQ2). Do process-derived 
features add predictive value beyond traditional metrics? (E2, RQ3).

Taken together, these evaluation steps provide a comprehensive analysis of feasibility, predictive performance, and feature contribution, ensuring that the proposed method is validated both structurally and empirically.